\section{Experimental Setup}
\subsection{Dataset}
\textbf{Sketches.} We built \CountSketches{} UI sketches, four for each of \CountCategories{} screen categories: login, registration, dashboard, profile, e-commerce, banking, booking, contact, settings and landing page. Each screen is specified as a component tree with varied layouts (navigation bars, multi-column rows, cards, image placeholders, radio and checkbox groups, selects, multiple buttons). It is rendered in a hand-drawn style with jittered, overshooting multi-segment strokes, slightly rotated and non-closed rectangles, varying pen width, three handwriting fonts and paper noise. Ground truth (type, label, normalised box and parent for \CountGtComponents{} components) comes from the specification. Each sketch also has a natural spoken description used in the generation step, plus the interactions and requirements that description implies. The sketches are \emph{synthetic}; Section~\ref{sec:limits} discusses what this implies.

\textbf{Voice commands.} We wrote \CountCommands{} commands, ten for each of add, remove, modify, move, resize, style, interaction and navigation. They include all seven example commands of the project specification (e.g., ``Make the login button blue'', ``Add a forgot password link below the login button''). Each command is tied to a sketch, labelled with its expected operation, target and parameters, and has a machine-checkable post-condition evaluated on the rendered UI. Each command was synthesised to speech with a text-to-speech model, rotating six voices, so that Whisper transcribes real audio.

\subsection{Baseline}
The baseline is the direct \emph{sketch$\rightarrow$AI$\rightarrow$UI} approach: one call to the same vision model (GPT-4.1-mini) receives the sketch and the spoken instruction and returns a complete React program. Modification sends the previous program and the command, and receives a regenerated program. The baseline prompt is a competent, fair prompt, and its output follows the same runtime contract, so the same renderer measures both systems.

\subsection{Metrics}
All metrics are computed by the same code for both systems. Measurement never uses the pipeline's internal \texttt{data-ngid} attributes; components are recovered from the rendered DOM by tags, ARIA roles, input types and visual cues.
\begin{itemize}
\item \textit{Component detection}: Hungarian matching~\cite{kuhn1955hungarian} of predicted to ground-truth boxes (IoU $\geq 0.3$). A match counts as a true positive when the types agree. Precision, recall and F1 are micro-averaged. It is measured on the analysis output (\texttt{cv}/\texttt{vision}/\texttt{hybrid}) and on the rendered DOM of each system.
\item \textit{Voice command accuracy}: Whisper word error rate (WER), plus the share of commands whose operation type, target and parameters are all interpreted correctly, from both the reference text and the Whisper transcript.
\item \textit{Visual similarity}: mean layout IoU between ground-truth components and their assigned rendered elements, and SSIM~\cite{wang2004ssim} between dilated Canny~\cite{canny1986edge} edge maps of the sketch and the screenshot.
\item \textit{Functional correctness}: the share of specified interactions that work when simulated in the browser. \textit{Requirement satisfaction}: the share of instruction-derived requirements (named components, input types, behaviours) that are met.
\item \textit{Refinement success}: the share of commands whose post-condition holds on the refined UI. \textit{Preservation}: the share of non-target components unchanged in type, text, colour, size and font.
\item \textit{Validation pass rate} before and after repair, and a \textit{fault-injection} study: deleted, hidden, zero-size and reordered nodes, broken events and injected runtime errors, scored by whether the validator detects each fault and whether repair removes it.
\item \textit{Generation time} per stage and end to end, and \textit{code quality}: parse success, lines of code, components, reuse ratio (component usages per definition), longest function and ESLint messages.
\end{itemize}

\subsection{Implementation}
The backend is Python (FastAPI, OpenCV, Playwright, SQLite for projects and versions), and the studio front end is React. The models are GPT-4.1-mini for vision and intent and \texttt{whisper-1} for speech, through the OpenAI API. The evaluation reran every stage in-process with cached raw model outputs, so metrics can be recomputed without further API calls. The full run made \ApiCalls{} model calls.
